%======================================================================
\section{Síntesis de resultados: de lo determinista a lo probabilístico}\label{sec:riesgo-sintesis}
%======================================================================
Las secciones anteriores siguen el orden en que se hicieron los experimentos. Esta sección los
reordena por peldaño de la escalera (figura~\ref{fig:riesgo-escalera}): primero lo que logra un
MPC determinista, sin y con pruebas de fallas, y qué tan robusto es; luego qué agrega cada manera
de incorporar el riesgo.

\subsection{Peldaño 0 sin fallas: el MPC determinista resuelve bien el problema que tiene}
\begin{itemize}
  \item \textbf{Lazo abierto.} N3 resuelve P0 en $\sim0.05$~s y N2 en cuatro QP de 10~ms, con la
        misma decisión; una cota inferior certificada la ubica a menos del 2.3\,\% y del 11.4\,\%
        del óptimo global en las dos consultas de prueba (parte~IV).
  \item \textbf{Lazo cerrado nominal} (50 realizaciones, sección~\ref{sec:desempeno}). N1, N2 y N3
        dan la misma trayectoria, se asientan en el 100\,\% de las realizaciones (5.5~s con
        40~kg/s, 13~s con la caja archivada) y la PCT no pasa de 1185~K. Su IAE es entre 6 y 29
        veces menor que el del PID y el del PI de la tesis.
  \item \textbf{Incertidumbre que no importa.} El error de pronóstico de la fuga mueve la PCT
        predicha $\sim0.025$~K y no cambia la decisión (F4, sección~\ref{sec:riesgo}). El desajuste
        ROM--HRM es un sesgo de $\sim0.2$~K que la corrección \emph{offset-free} elimina
        (sección~\ref{sec:f45}). Un pronóstico de Kalman o el caudal medido del ACS cambian el caudal
        hasta 4~kg/s, pero no la PCT ($0.00\pm0.00$~K, sección~\ref{sec:riesgo-donde}).
\end{itemize}

\subsection{Peldaño 0 con pruebas de fallas: cuán robusto es el MPC determinista}
El MPC determinista no conoce las fallas, pero se puede \emph{probar} contra ellas. Se hicieron
cuatro pruebas, de la más dura a la más realista (tabla~\ref{tab:sintesis-robustez}).

\begin{table}[H]\centering\scriptsize\setlength{\tabcolsep}{3pt}
\begin{tabularx}{\linewidth}{@{}>{\raggedright\arraybackslash}p{30mm} >{\raggedright\arraybackslash}p{36mm} Y >{\raggedright\arraybackslash}p{28mm}@{}}\toprule
Prueba & Qué falla & N1 / N2 / N3 & Comparación\\\midrule
Familia extrema, caja archivada (12.9~kg/s; sec.~\ref{sec:extremos}) & fuga 15--25~kg/s, ACS al 60--90\,\%, fallas de 30--90~s & 9/14 y 22/50 llegan a 1220~K; el MPC satura siempre y da lo mismo que el caudal máximo & sin ACS: 14/14\\
Familia extrema, actuador físico (sec.~\ref{sec:extremos}, \ref{sec:desempeno}) & ídem & 0/14 y 0/50; PCT máxima 1189--1207~K, margen mínimo 13~K & PI del DET: 0/14; el MPC usa 25\,\% menos masa\\
Mapa de capacidad (sec.~\ref{sec:riesgo-donde}) & capacidad 15--70~kg/s, familia extrema & con $\le20$~kg/s, 3--6 de 14 cruzan; con $\ge40$~kg/s, ninguna & PI: 6--10 de 14\\
Fallas del DET en el lazo, 40~kg/s (sec.~\ref{sec:riesgo-fallas}) & 5 componentes Markov, $\lambda\times10^5$, $\mu\times10^4$, $\ge1$ falla por realización & 0/50; PCT pico media 1189~K; DSM mínimo 54.9~K & PID 1/50, PI 7/50\\
Riesgo integrado (Rao--Blackwell, sec.~\ref{sec:riesgo-rb}) & cualquier instante y duración de la primera falla & $P(\text{cruce}\mid\text{falla})=0.031$; $\approx1.0$\,\% por transitorio & PID 0.033, PI 0.119\\
Sensibilidad a la duración (sec.~\ref{sec:riesgo-aplicabilidad}) & media 15, 46 o 150~s & 0.002, 0.031 y 0.32 & ---\\\bottomrule
\end{tabularx}
\caption{Robustez del MPC determinista frente a fallas del ACS.}
\label{tab:sintesis-robustez}
\end{table}

\paragraph{Qué significa «robusto» aquí.} Conviene separar tres afirmaciones.
\begin{enumerate}
  \item \textbf{Robustez empírica}: en las familias probadas con capacidad suficiente, ninguna
        realización de N1--N3 cruza 1220~K, y los controladores de la tesis sí lo hacen. Es una
        prueba de estrés: dice que no se vio un cruce, no cuán probable es.
  \item \textbf{Riesgo residual}: medido con Rao--Blackwell, N2 cruza en el 3.1\,\% de las fallas
        del caso base. Ese número no lo eligió nadie: lo fija la consigna.
  \item \textbf{Garantía}: el MPC determinista no tiene ninguna. Su límite $U_y$ nunca se activa
        (proposición~\ref{prop:blanda-dura}), y su predicción no contiene la falla.
\end{enumerate}

\paragraph{De dónde viene la robustez.} No de anticipar la falla, sino de \emph{seguir bien la
consigna}. Con la PCT en $y_{\rm ref}\approx1180$~K hay 40~K hasta 1220~K, y la planta tarda
$\tau\approx160$~s en recorrerlos sin ACS; sólo cruzan las fallas más largas que eso,
$S(160)=0.031$ (figura~\ref{fig:riesgo-tutorial}b). El PI del DET tiene una zona muerta y recibe
las fallas con la PCT alta: su $\tau$ baja hasta 45~s y su riesgo es casi cuatro veces mayor. La
consigna actúa como un \emph{margen implícito}. La robustez del MPC determinista es, por lo tanto,
la de su consigna más la precisión con que la sigue.

\paragraph{Dónde se rompe.} (i) Si la capacidad del actuador no alcanza a la fuga (caja
archivada), ningún controlador ni capa de riesgo evita el cruce: es un problema de diseño del ACS,
no de control. (ii) Si las fallas duran tanto como $\tau$ (media de 150~s), el riesgo sube a 0.32
y, de nuevo, ninguna capa lo baja a un nivel aceptable. (iii) Si las fallas se recuperan en
decenas de segundos (caso base), el margen implícito deja un riesgo de 3\,\% que \emph{sí} se puede
reducir pre-enfriando, y ése es el único régimen en que los peldaños siguientes cambian algo.

\subsection{Peldaño 1: el margen determinista no agrega protección}
El DSM de la tesis no cambia la trayectoria de N2 con ningún $\lambda_{\rm dsm}$ entre $10^2$ y
$10^6$, ningún límite (1244 o 1220~K) ni ningún umbral probado, con y sin fallas: la PCT pico
cambia menos de 0.002~K y la trayectoria con fallas del DET es la misma, con algo más de
movimiento del actuador (secciones~\ref{sec:dsm-lambda} y~\ref{sec:riesgo-fallas}). La razón es
estructural: el DSM se calcula sobre la predicción nominal, que no contiene la falla, y su umbral
está 0.15~K sobre la consigna que N2 ya sigue. En la tesis se activaba por los errores del
optimizador (LJ) y de la predicción, no por el riesgo. Un DSM con umbral bajo la consigna sería
una consigna más baja, fija, sin calibrar en probabilidad.

\subsection{Peldaño 2: las métricas probabilísticas}
\paragraph{Riesgo de pronóstico (F4).} La SAA de dos etapas con 59 trayectorias de fuga, sobre
N3, corre en $\sim2$~s y da la misma primera acción que el problema determinista: la incertidumbre
de la fuga es unas 100 veces menor que la autoridad del actuador.

\paragraph{Riesgo de fallas del ACS (F5b--F8).} La tabla~\ref{tab:sintesis-prob} reúne los
resultados con 50 realizaciones a 40~kg/s, $\mu\times10^4$ y $R=0.1$ donde corresponde.

\begin{table}[H]\centering\scriptsize\setlength{\tabcolsep}{3pt}
\begin{tabularx}{\linewidth}{@{}>{\raggedright\arraybackslash}p{38mm} r r r >{\raggedright\arraybackslash}p{18mm} Y@{}}\toprule
Formulación (peldaño) & $P(\text{cruce}\mid\text{falla})$ & Pre-enfr.\ [K] & Masa & Fallas en el lazo & Garantía / teorema\\\midrule
N2 (0) & 0.031 & 0.0 & --- & 0/50 & ninguna; margen implícito\\
N2 + DSM, $\lambda=10^4$ (1) & 0.031 & 0.0 & --- & 0/50 & ninguna\\
Pérdida esperada $\rho=100$ (2) & 0.026 & 1.8 & +4\,\% & 0/50 & equilibrio de costos; SAA\\
Cuantil $\alpha=0.05$ (2) & 0.027 & 1.1 & +2\,\% & --- & $P\le\alpha$; lema del cuantil\\
CVaR$_{0.90}$ (2) & 0.021 & 3.6 & +8\,\% & --- & severidad de la cola; R--U\\
Cuantil $\alpha=0.02$ (2) & 0.015 & 7.4 & +17\,\% & 0/50 & $P\le\alpha$, verificada\\
CVaR$_{0.95}$ (2) & 0.014 & 8.6 & +20\,\% & 0/50 & severidad; R--U\\
Pérdida esperada $\rho=1000$ (2) & 0.014 & 9.4 & +21\,\% & --- & equilibrio de costos; SAA\\
Wilks 59 (2) & 0.011 & 13.0 & +29\,\% & 0/50 & 95/95; regla de Wilks\\
Cuantil $\alpha=0.01$ (2) & 0.011 & 12.9 & +32\,\% & --- & $P\le\alpha$, en el límite\\
Un escenario aprendido, $\rho=100$ (3) & 0.026 & 1.8 & --- & --- & = SAA; Homem-de-Mello\\
Un escenario aprendido, $\rho=1000$ (3) & 0.015 & 8.2 & --- & --- & 99.7\,\% del VSS\\\bottomrule
\end{tabularx}
\caption{Las formulaciones probabilísticas ordenadas por riesgo (tablas~\ref{tab:rb},
\ref{tab:det-tesis} y~\ref{tab:f9-lazo}). «Fallas en el lazo»: cruces de 1220~K con las fallas
del DET aplicadas a la planta (campaña F6 \code{det}); «---»: no corrido en esa campaña.}
\label{tab:sintesis-prob}
\end{table}

\paragraph{Lo que muestran, en orden.}
\begin{enumerate}
  \item \textbf{La métrica probabilística sí cambia la decisión}, y en la cantidad que pide su
        parámetro: la cota con cuantil cumple $P\le\alpha$, verificado sobre la planta. El
        mecanismo es único: pre-enfriar antes de la falla. Durante la falla todas las formulaciones
        se calientan igual (sección~\ref{sec:trayectorias}).
  \item \textbf{Todas recorren una sola frontera.} Como el pico es aditivo en la duración, el
        cuantil, el CVaR, Wilks y el escenario único se reducen a un margen, y la pérdida esperada
        cae sobre la misma curva (proposición~\ref{prop:margen}). Ninguna compra más protección por
        kelvin de pre-enfriamiento.
  \item \textbf{El precio es explícito}: bajar el riesgo a la mitad cuesta 7.4~K bajo la consigna
        y un 17\,\% más de plomo; a un tercio, 12.9~K y 32\,\%. Con las fallas en el lazo, el IAE
        de la tesis sube (de 475 a 1689~K\,s con $\alpha=0.02$) porque cuenta el margen como error,
        y el DSM mínimo sube de 54.9 a 59--61~K.
  \item \textbf{Las extensiones no cambian la conclusión}: con fallas parciales el riesgo es $p_0$
        veces el de las totales y la cota con $\eta$ es la de falla total con $\alpha/p_0$; la
        capacidad del ACS entre 30 y 50~kg/s casi no influye; con $R\ge0.03$ desaparece la oscilación
        de las cotas sin cambiar el riesgo.
  \item \textbf{La utilidad depende de un dato físico}: el tiempo de recuperación automática del
        ACS frente a $\tau\approx165$~s (sección~\ref{sec:riesgo-aplicabilidad}).
\end{enumerate}

\subsection{Peldaño 3: el escenario único de Rozas}
El teorema de Homem-de-Mello et al.\ se cumple en el SMDFR para la pérdida esperada y, escritas
por paso, para las fallas parciales; la unicidad la da $R$. El escenario exacto reproduce la
decisión de la SAA pero no es un evento físico; una sola falla bien elegida, de 170--190~s (cuatro
veces la media), recupera el 94--99.5\,\% del VSS, y la duración media no recupera nada. Aprendido
según el contexto, con el método del artículo, captura el 89--97\,\% del VSS fuera de muestra y
en lazo cerrado reproduce el riesgo de la SAA, si el contexto sólo usa variables que la decisión no
mueve. En el SMDFR actual la ganancia es de holgura de cómputo (0.07 frente a 0.42~s en el peor
lazo), no de decisión (sección~\ref{sec:rozas}).

\subsection{La recomendación, con su razón}
\begin{center}\small
\begin{tabularx}{\linewidth}{@{}>{\raggedright\arraybackslash}p{40mm} Y@{}}\toprule
Elección & Razón\\\midrule
N2 + corrección de sesgo & misma decisión que N3, certificada; QP convexo de $\sim30$~ms\\
Rampa dura, límite $U_y$ blando & la rampa siempre es factible; el límite nunca se activa\\
$R=0.1$ & suaviza el caudal y elimina la oscilación de la capa de riesgo sin costo\\
Cota con cuantil, $\alpha$ a elegir con el equipo & la métrica más simple y con garantía verificada; las demás no son más eficientes\\
Evaluación por Rao--Blackwell sobre la planta & 150 veces menos realizaciones que contar cruces\\
SAA o escenario único & sólo si se prefiere penalizar el exceso; dan la misma decisión\\\bottomrule
\end{tabularx}
\end{center}

\takeaway{\textbf{En orden.} El MPC determinista (N1--N3) resuelve bien su problema y, con
capacidad suficiente, no cruza 1220~K en ninguna prueba de fallas; su robustez viene de seguir
con precisión una consigna que deja 40~K de margen, y su riesgo residual es del 3\,\% por falla.
El margen determinista de la tesis no agrega nada. Una métrica probabilística convierte ese margen
implícito en uno calibrado: la cota con cuantil cumple $P\le\alpha$ pagando pre-enfriamiento, y
todas las métricas recorren la misma frontera. El escenario único de Rozas reproduce la decisión
estocástica con un solo QP. Lo que ninguna de estas formulaciones hace es adaptar decisiones
futuras al modo del ACS: eso es el MPC multietapa.}

%======================================================================
\section{Problemas abiertos y prometedores}\label{sec:riesgo-futuro}
%======================================================================
\subsection{El MPC multietapa (vía V8 de la ruta)}
\paragraph{Qué falta en las formulaciones actuales.} Todas las capas de riesgo son de \emph{dos
etapas}: un solo plan $u\in\mathbb R^{26}$ para todos los escenarios, y una segunda etapa que sólo
\emph{mide} el exceso sobre 1220~K. En un MPC multietapa, en cambio, las decisiones futuras
dependen de lo que se haya observado: el plan se ramifica según el modo del ACS. Con
$\sigma_n\in\{0,1\}$ (ACS operativo o fallado) en cada nodo $n$ de un árbol, con padre $a(n)$ y
probabilidad $\pi_n$:
\begin{align*}
\min_{\{u_n\}}\quad & \sum_n\pi_n\Big[(\hat y_n-y_{\rm ref})^2+R\,(u_n-u_{a(n)})^2+\rho\,(\hat y_n-1220)_+\Big]\\
\text{s.a.}\quad & x_n=\Phi\big(x_{a(n)},\,(1-\sigma_n)\,u_n,\,\hat d_n\big),\qquad \hat y_n=h(x_n)+d,\\
& 0\le u_n\le\mdot_{\max},\qquad |u_n-u_{a(n)}|\le\Delta\mdot .
\end{align*}
La \emph{no anticipación} está en la indexación: dos escenarios que comparten la historia hasta un
nodo comparten la decisión de ese nodo. La versión de dos etapas actual es el caso en que todos los
nodos de una misma etapa comparten $u$. En la práctica se ramifica sólo en los primeros pasos
(\emph{horizonte robusto}) y el resto del árbol se deja sin ramificar, como en el MPC multietapa
no lineal de Lucia et al.~\cite{lucia2013}; ver también la revisión de MPC estocástico de
Mesbah~\cite{mesbah2016} y el texto de Birge y Louveaux~\cite{birge2011}.

\paragraph{Por qué hoy rendiría poco.} Hay tres razones, todas medidas en este documento.
(i) Durante la falla el caudal es cero: no hay decisión que ramificar. (ii) Tras la reparación,
el MPC re-optimiza cada 0.5~s, y eso ya es un recurso. (iii) El horizonte es de 13~s y las fallas
duran decenas de segundos: dentro del horizonte el árbol casi nunca cambia de modo. Lo único que se
decide es el margen \emph{antes} de la próxima falla, que es lo que la segunda etapa ya captura.

\paragraph{Cuándo rendiría.} El multietapa ganaría valor si aparece una decisión que dependa del
modo del ACS:
\begin{enumerate}
  \item \textbf{Varias protecciones} (vía V9): con un segundo sistema activo o el tapón de sal
        pasivo, durante una falla del ACS sí hay que decidir qué activar y cuándo. A$\to$B y
        B$\to$A pueden dejar riesgos distintos.
  \item \textbf{Tiempos con memoria}: retardos de detección y de actuación, o reparaciones de
        duración casi fija. Entonces la edad de la falla predice cuánto falta, y una política que
        la observe puede decidir mejor. La mezcla del DET tiene un poco de memoria
        (ejemplo~\ref{ej:memoria}), pero durante la falla no hay nada que decidir.
  \item \textbf{Un horizonte de riesgo largo}: para que el árbol vea la reparación, el horizonte
        debe cubrir $q_{1-\alpha}(D)\approx180$--220~s. Con pasos de 0.5~s son 360--440 pasos por
        rama; haría falta un bloqueo de movimientos (vía V7) o una malla gruesa para la parte de
        riesgo.
\end{enumerate}

\paragraph{Cómo medir cuánto se pierde con dos etapas.} La brecha entre el MPC de dos etapas
(nivel B de la sección~\ref{sec:rozas-alcance}) y la política multietapa óptima (nivel A) se puede
medir con programación dinámica sobre un modelo reducido. El estado es el margen $m$ hasta
1220~K, el modo $\sigma$, la edad $a$ de la falla y, si hay varias protecciones, la etapa de
protección; la ecuación de Bellman es
\[
V(m,\sigma,a)=\min_{u\in\mathcal U(\sigma)}\Big\{c(m,u)+\mathbb E\big[V(m',\sigma',a')\mid m,\sigma,a,u\big]\Big\},
\]
con $m'$ dado por la dinámica reducida (enfriamiento con $u$ si $\sigma=0$; calentamiento de
$\approx0.3$~K/s si $\sigma=1$) y las transiciones de $(\sigma,a)$ dadas por las tasas del DET. Con
un estado de tres o cuatro dimensiones, el problema se resuelve en una malla. Si la política de
dos etapas queda cerca de $V$, el multietapa no vale la pena para el SMDFR actual; si no, la
brecha dice cuánto se gana y en qué estados.

\subsection{El escenario único más allá de dos etapas}
El teorema de Homem-de-Mello et al.\ trata problemas de dos etapas: cierra la brecha entre un
escenario y el MPC de dos etapas, no la brecha con la política multietapa
(sección~\ref{sec:rozas-alcance}). Una extensión natural, que queda como \emph{pregunta de
investigación} y no como resultado del artículo, es aplicarlo etapa por etapa. En los métodos de
descomposición anidada, como la programación dual dinámica estocástica (SDDP~\cite{pereira1991}),
cada etapa resuelve un problema de dos etapas cuya segunda etapa es una función de costo futuro
convexa y poliédrica, construida con cortes. Si cada uno de esos problemas cumple las condiciones
del teorema~1 (recurso fijo, costos fijos, duales acotados), existe un escenario por etapa que
reproduce su decisión, y la estrategia de Rozas pasaría a aprender una \emph{secuencia} de
escenarios. Habría que comprobar tres cosas: (i) que el costo futuro conserve la estructura que el
teorema pide; (ii) que el lema~1 se cumpla etapa por etapa cuando la falla multiplica a la decisión
($T$ aleatoria); y (iii) que el contexto siga siendo exógeno, que es lo más delicado: en un
problema multietapa el estado que pasa de una etapa a la siguiente depende de la decisión, el mismo
fenómeno que hizo oscilar el lazo en E5 (sección~\ref{sec:rozas-lazo}).

\subsection{Otros problemas prometedores}
\begin{itemize}
  \item \textbf{Robustez frente a la distribución de la duración.} El error de especificación de
        $\mu$ es asimétrico: si las fallas duran más que las de diseño, la cota con $\alpha=0.02$
        deja 0.26 (sección~\ref{sec:riesgo-aplicabilidad}). Una cota \emph{distribucionalmente
        robusta}~\cite{delage2010}, que exija $P\le\alpha$ para todas las duraciones medias
        plausibles, se reduce aquí a diseñar con la recuperación más lenta del conjunto, porque
        $S$ crece con la duración media; lo que falta es el conjunto, que debe salir del equipo.
  \item \textbf{Riesgo por transitorio y varias fallas.} La cota controla una falla que empieza en
        la consulta. Una restricción conjunta sobre todo el transitorio, o con dos fallas seguidas,
        pediría un horizonte de riesgo largo y, probablemente, el árbol del multietapa.
  \item \textbf{Incertidumbre del modelo de calentamiento.} El supuesto A1 es conservador en
        $\sim10$\,\% en la réplica del HRM; falta validarlo nodo a nodo durante una falla contra
        Simulink. Si el error resultara del lado inseguro, se agregaría como back-off gaussiano
        (proposición~\ref{prop:cc}) con la varianza del error ROM--HRM durante la falla.
  \item \textbf{Escenario aprendido con contexto exógeno.} La estrategia de Rozas rendiría más con
        una segunda etapa cara (varias fallas, fallas parciales por paso) y con contexto que la
        decisión no mueve: la etapa del sistema de protección, la edad de la falla o la fuga
        medida.
\end{itemize}

\begin{table}[H]\centering\scriptsize\setlength{\tabcolsep}{3pt}
\begin{tabularx}{\linewidth}{@{}>{\raggedright\arraybackslash}p{36mm} Y >{\raggedright\arraybackslash}p{36mm} c@{}}\toprule
Problema & Por qué promete & Qué haría falta & Prioridad\\\midrule
Tiempo de recuperación del ACS & decide el régimen en que la capa de riesgo sirve & dato del equipo (trenes, arranque automático) & 1\\
Referencia por programación dinámica & mide la brecha dos etapas--multietapa antes de construir el árbol & modelo reducido $(m,\sigma,a)$; días de trabajo & 2\\
MPC multietapa con varias protecciones (V8--V9) & primera fuente real de decisiones dependientes del modo & modelo de la segunda protección; árbol con horizonte robusto; bloqueo de movimientos & 3\\
Cota robusta en $\mu$ & la garantía actual se pierde si las fallas duran más & conjunto de duraciones plausibles & 3\\
Escenario único multietapa (Rozas + SDDP) & extendería la estrategia a la política multietapa & comprobar el teorema por etapa; contexto exógeno & 4\\
Validar A1 nodo a nodo & cierra la última brecha de modelo del resultado de riesgo & MATLAB (leo\_asus) & 4\\\bottomrule
\end{tabularx}
\caption{Problemas abiertos de la capa de riesgo, por prioridad.}
\label{tab:riesgo-futuro}
\end{table}

\takeaway{\textbf{Hacia dónde seguir.} Con las fallas del DET y un solo sistema de protección,
las formulaciones de dos etapas ya capturan la única decisión que importa: cuánto margen tener
antes de la falla. El MPC multietapa, y con él una versión multietapa del escenario único de Rozas,
valdría la pena cuando aparezcan decisiones durante la falla (varias protecciones), tiempos con
memoria u horizontes de riesgo largos. Antes de construirlo conviene medir la brecha con una
referencia de programación dinámica sobre un modelo reducido.}
